% Quantitative choice of atom budgets and uniform scale thresholds.
\subsection{Choosing the local paid budgets}
\label{sec:ex-lattice-local-paid-budgets}

The one-step insertion uses two kinds of local input: a finite list of
coefficient-recovery probes, and positive budgets for the basis, regular,
and forced atoms. The first choice is elementary once the probe list is
finite. The eventual threshold can then be made monotone in the scale and
uniform over all admissible adjacent levels. This section proves those two
selection steps and then matches their budgets to the actual basis,
regular and forced insertions. Those estimates supply the one-step
recurrence of Subsection~\ref{sec:ex-lattice-trace-price-recurrence}.

Fix a finite set \(\mathcal T\) of probes and nonnegative constants
\(O,H,W,G_d,a,c\), where \(O\) is the motion-average norm factor, \(H\) the local
readout factor, \(W\) the test factor, \(G_d\) the polarization factor,
\(a\) the block weight, and \(c\) the tail-decay factor. For each probe
\(\tau\), let \(\pi_\tau\geq0\) be its variation price and define its
local coefficient
\[
 q_\tau=4aG_d\pi_\tau+2a.
\]
Given a target price \(y>0\), the required atom budget inequalities are
\begin{align}
 q_\tau v&\leq u &&(\tau\in\mathcal T), \label{eq:ex-lattice-atom-cycle}\\
 O\bigl(Wu+2Hav\bigr)&\leq y, \label{eq:ex-lattice-atom-regular}\\
 4OHav&\leq y, \label{eq:ex-lattice-atom-large}\\
 O\bigl(Wv+2Havc\bigr)&\leq y. \label{eq:ex-lattice-atom-main}
\end{align}
The two outputs \(u,v\) are the probe and summability budgets. In the
Wilson application, at coefficient grade \(d\) and block side \(L\), the
constants are specialized as follows: \(\mathcal T\) is the finite
coefficient-recovery probe family \(\mathcal Q_d\) of
Subsection~\ref{sec:ex-lattice-trace-price-recurrence} and \(W=T_d\) its
reconstruction constant; \(G_d\) is the polarization constant \(M_d\) and
\(\pi_\tau=\pi(\tau)\geq1\) the probe price defined there; \(O\) is the
orbit-average cost and \(H\) the readout price of the grade-\(d\)
coefficient frame; \(a=L^d\); and \(c\) is the credibility of the fine profile,
that is, its uniform mass bound for a normalized frame-column insertion,
fixed by its support exponent, order and power. This is the constant called $h$ in
Subsection~\ref{sec:ex-lattice-trace-eve-producer}. With
these values \(q_\tau\) bounds the local variation coefficient,
while the other three inequalities pay respectively for regular terms,
large-support terms, and the primary atom. The same number \(v\) is later
used as the summability budget in all of the local atom estimates of one
transition; the budgets are not chosen separately for each estimate.

\begin{lemma}[Positive local budget choice]
\label{lem:ex-lattice-local-budget-choice}
For every finite probe family and every \(y>0\), there are \(u,v>0\)
satisfying \eqref{eq:ex-lattice-atom-cycle}--\eqref{eq:ex-lattice-atom-main}.
\end{lemma}
\begin{proof}
Set
\[
 K=\sum_{\tau\in\mathcal T}q_\tau,\qquad
 B=\max\{1,\ O[W(1+K)+2Ha],\ 4OHa,\ O[W+2Hac]\}.
\]
These are finite nonnegative numbers. Choose
\[
 v=\frac{y}{2B},\qquad u=(1+K)v.
\]
Both are strictly positive. Since \(q_\tau\leq K\leq1+K\),
\(q_\tau v\leq u\). The left side of the regular-term inequality is
\(vO[W(1+K)+2Ha]\leq vB=y/2\). The large-support and primary left
sides are respectively \(4OHav\leq Bv=y/2\) and
\(vO[W+2Hac]\leq Bv=y/2\). Each is therefore at most \(y\), as
required.
\end{proof}

We next record the uniform threshold selection used to turn fixed-parameter
estimates into a single scale profile. Here \(j\) is a stage index of the
selected stage data, with profile orders \(\operatorname{order}_i\),
powers \(\operatorname{power}_j\) and first stage \(i_0\). At a stage
\(j\), the admissible pairs of an adjacent stage \(i\) and a coefficient
grade \(d\) form the set
\[
 \mathcal I_j=\bigl\{(i,d):\ i\leq j\leq i+1,\ i_0\leq i,\
   d+1\leq\operatorname{order}_i,\ d\leq\operatorname{power}_j,\
   d+2\leq j\bigr\}.
\]
The adjacent stage is \(j\) itself or its immediate predecessor, so
there are at most two, and \(d+2\leq j\) confines the grades to the finite
set \(\{0,\ldots,j-1\}\). Thus \(\mathcal I_j\) is finite, and only
finitely many eventual requirements
occur at each fixed \(j\).

\begin{lemma}[Monotone threshold envelope]
\label{lem:ex-lattice-monotone-threshold}
Let \(b_j\in\mathbb R\) be any prescribed lower threshold. For each
\(j\), let \(S_j\) be a finite set and, for each \(s\in S_j\), let
\(E_{j,s}\in\mathbb R\) have the property that a desired conclusion
\(\mathcal P_{j,s}(x)\) holds whenever \(x\geq E_{j,s}\). Then there is a
nondecreasing sequence \(B_j\) such that \(B_j\geq b_j\) and, for every
\(s\in S_j\), \(\mathcal P_{j,s}(x)\) holds whenever \(x\geq B_j\).
\end{lemma}
\begin{proof}
Define
\[
 e_j=\max\bigl(\{b_j\}\cup\{E_{j,s}:s\in S_j\}\bigr),
 \qquad B_j=\max_{0\leq r\leq j}e_r.
\]
The maxima exist because the sets are finite and nonempty after adjoining
\(b_j\). The prefix maximum \(B_j\) is nondecreasing and dominates
\(b_j\). If \(s\in S_j\), then \(E_{j,s}\leq e_j\leq B_j\), so every
\(x\geq B_j\) satisfies the stipulated conclusion.
\end{proof}

Apply this lemma at each stage \(j\) with \(S_j=\mathcal I_j\). Let a
positive target price \(y_{j,i}\) be prescribed for every pair of stages.
For each fixed \((i,d)\in\mathcal I_j\), the
atom-budget lemma, with \(y=y_{j,i}\), first supplies positive budgets
\(u,v\); for these fixed budgets, the
finite-volume local estimates then supply an eventual threshold
\(E_{j,(i,d)}\) in the scale variable \(x\) for the
corresponding primary, regular, and forced atom conclusions. These local
estimates require, in addition, that the smallness parameter
\(\varepsilon\) attached to the transition \((j,i)\) is at most \(1/4\). Taking the
finite maximum and then the prefix maximum gives one monotone threshold
profile that works simultaneously for every adjacent transition and every
permitted coefficient grade. This is the uniformization step needed before
the scale induction. The atom estimates to which this uniformization applies are established
and matched below.


\subsubsection{The three atom estimates used in the recurrence}
\label{subsec:trace-paid-atom-suppliers}
We verify the estimates for the selected transition, with the same
budget choices $u,v$ as above. Put $\rho=L^{-2/5}$ and
$\chi=\sqrt{H(x)}$. The selected response package is
Theorem~\ref{thm:ex-lattice-eve-producer}, with its probe and cloud
bounds proved in Subsection~\ref{subsec:trace-eve-probe-cloud}.
In particular, every cloud of fine mass $m$ has charge at most
$vm/\chi$ on a common tail. The two real/imaginary channel charges
may be equal; they are nevertheless counted twice.

\subparagraph{Regular atoms.}
For a regular source of mass $m$, use the actual small- and large-support
constructions in Subsection~\ref{subsec:trace-regular-chart-construction}.
In the small case the local probe charge is at most
$q_\tau vm/\chi\le um/\chi$ for every recovery probe.
Coefficient recovery costs $W$, and the two cloud readouts cost at most
$2Havm/\chi$. Including the motion factor $O$ gives
\[
 \frac{m}{\chi}O(Wu+2Hav)\le \frac{ym}{\chi}.
\]
In the large case each of the two chart terms and each of the two
cloud terms is at most $vm/\chi$. Hence the total is
$4OHavm/\chi\le ym/\chi$.
Proposition~\ref{prop:ex-lattice-regular-source-atom} gives the
associated presentations, their exact pushed density, and remainder
cost at most $\rho m$ with the prescribed coarse forced normalization.
Its derivative and vanishing-jet assumptions hold for the regular
atoms of the effective-representation certificate; they are not inferred
from a pointwise bound. Thus the cost pair is $(ym/\chi,\rho m)$.

\subparagraph{Basis atoms.}
Let $e_\alpha$ be a nonconstant normalized basis column, and insert its
actual local section with the fine threshold. The cloud mass in its
primary correction is bounded by the fine-profile credibility $c$. The primary expansion and coefficient extraction
are those of Subsection~\ref{sec:primary-calibration-producer}.
Choose the response package at price $v$. For a recovery probe $p$,
compare the actual principal row first to its Gaussian reference, then
to its cropped deterministic prediction, then to its full polynomial
forecast. The first two comparisons, in the two real/imaginary channels,
have ordered derivative error at most $2(e|u(x)|+\delta_x)$:
the insertion bound supplies $e|u(x)|$ in each channel, and the cap
bound supplies $\delta_x$ for the fluctuation-to-prediction comparison.
Polarization and the derivative prices of the probe curve cost
$\pi(p)\Lambda$; multiplication by the physical factor costs $L^d$.
The cropped-to-full comparison costs the additional $\delta_x$ by
\eqref{eq:trace-eve-probe-envelope}. Thus the actual principal probe
error is bounded by
\[
 \pi(p)\Lambda L^d\,2(e|u(x)|+\delta_x)+\delta_x
       \le C_x\le v/\chi,
\]
exactly the charged expression of
Theorem~\ref{thm:ex-lattice-eve-producer}. The Gaussian-to-prediction
comparison is the normalized derivative comparison of
Subsection~\ref{sec:primary-calibration-producer}; the cap pays its
Gaussian moment before any probe or column is chosen. Its full predicted coefficient is the leading map of that
same section, with both the amplitude-degree and spatial-grade factors
retained. The other rows have summed readout charge at most
$Ha(\mathcal M+\mathcal M)\le2Havc/\chi$ by the cloud bound.
The coefficient-extraction identity is linear; the finite recovery
map costs $W$, and its motionwise norm costs $O$. The error from the
prescribed leading coefficient is therefore at most
\[
 \frac1\chi O(Wv+2Havc)\le \frac y\chi.
\]
The graded principal remainder and the local splitting are proved in
Subsection~\ref{sec:primary-calibration-producer}: for a column of spatial
grade $g_\alpha$ they cost at most $P L^{d-g_\alpha}$, with the fixed
frame ledger constant $P$. The splitting preserves the exact density
by \eqref{eq:local-atom-identity}; it subtracts precisely the extracted
local section and puts the residual in the regular and forced lists.
It also preserves order-zero format, since the residual has zero
identity value and zero nonwrapping test jets. Covariance of extraction
under each lattice motion gives presentations of the motion of one
coefficient, rather than unrelated coefficients for different motions.
The unit column has its exact constant transport and no error.
By finite linearity, a coefficient $z$ consequently has error at most
$(y/\chi)\|z^\circ\|_1$ and remainder at most
$P\|\mathcal V_Lz^\circ\|_1$, where
$(\mathcal V_Lz)_\alpha=L^{d-g_\alpha}z_\alpha$.

\subparagraph{Forced atoms.}
Let a forced source have fine mass $m\ge0$, and set
$g=\sqrt{\operatorname{forcedAllowance}_{\rm fine}}\in(0,1]$.
We use the localized expansion with its marked profile and
address-background bounds from Subsection~\ref{sec:leaf-transfer}, with
the small bills funded as in Subsection~\ref{sec:projection-payment-producer}.
These are the hypotheses of the projection calculation for the actual
transition. For a forced insertion the principal cluster is zero: its flag is
incompatible with the healthy principal term. In the projection
calculation of Subsection~\ref{sec:leaf-transfer} we may therefore set
its principal price $V$ to zero. The two remaining bills and the
hit-width condition are paid in
Subsection~\ref{sec:projection-payment-producer}. With the notation
there, $e_0=(\rho/H_d)m$, so the motionwise coefficient cost and
the normalized coarse remainder cost are both at most
$H_de_0=\rho m$. The exact density and order-zero format are preserved
by that same splitting calculation. The flagged profile identity there
is required; equality of scalar densities alone would not justify it.

Lemma~\ref{lem:ex-lattice-medal-gas-price} gives
$\rho g\le y/\chi$ eventually. Consequently
\[
 \rho m\le\frac y\chi\frac m g,
 \qquad \rho m\le\rho\frac m g.
\]
These are the required forced cost pair, including $m=0$.
The first inequality uses the fine allowance in $g$; the remainder
on the left already uses the coarse forced normalization supplied by
the projection. No identification of those two allowances is made.

All statements concern the actual two effective representations of one
transition. The finite basis sum and the absolutely summable regular
and forced lists are combined in the density-summation argument of
Subsection~\ref{sec:ex-lattice-trace-price-recurrence}; the bounded-density
and summable-activity estimates there justify the interchange with
integration. At a fixed stage and grade there are only finitely many
probe, basis, channel and budget requirements. Their thresholds can
therefore be included in the finite maximum above, before the period,
volume or source is chosen. This is the promised source of the eventual
local estimates, not an additional choice of an atom-dependent threshold.
